\documentclass[11pt]{article}
\usepackage{amsmath,amssymb}
\usepackage[hidelinks]{hyperref}
\title{Rewinding a protected FIFO cut after supplier migration}
\author{}\date{}
\begin{document}\maketitle
\paragraph{Scope.}
This is an integer event model for the sender prefix check of the
preceding broker binding replay, not a new native integration. A row
has an immutable sequential rank, sender \(s\), destination \(e\),
payload size and (optionally) one named required receiver at position
\(k\) in its ordered recipient list. The operations are append, lawful
withdrawal of an unexecuted row, growth by appending a new receiver,
and atomic movement of a required row to the tail of a different sender's
queue. Positions of existing receivers therefore do not change.
Capacities are fixed effective sender allowances after any retained
old debit. Debit expiry, actual dequeue and charging, physical action,
authorization provenance and RX/path service are not modeled here.

\paragraph{Representation and invariants.}
For live rows define \(N_{s,e}(j)\) as the total bytes from \(s\)
to \(e\) at ranks strictly earlier than \(j\), and let \(m_e\)
be the number of bound receivers. For a protected row \(j\) of size \(b_j\),
the no-expiry projected debit through its named receiver is
\[
       H_s(j) = \sum_e m_e N_{s,e}(j) + k_j b_j .
\]
With nonnegative bytes and \(1\le k_j\le m_{e(j)}\), the full
prefix through the latest protected row dominates every earlier
required TX prefix for that sender. Retain a slack scalar for that
latest row, the weighted total for all pending rows, a heap of
protected ranks, and a sparse partial-sum tree per live sender-endpoint
pair, with an inverse map from endpoints to live senders. The tree is a
standard \href{https://doi.org/10.1002/spe.4380240306}{Fenwick tree}
on a fixed rank universe. It records logical bytes, not a prediction of
which attempts will actually charge.

Appending a protected row evaluates it against the current total plus
its named-receiver cost and makes it latest only if admitted.
Withdrawal of a row before the latest cut releases precisely its full
weighted term; removal strictly after it has no effect on certified
slack. Removing the latest row exposes a smaller cut: recompute at
the new rank, not at the old rank. A growth at \(e\) changes the
total of each incident sender by its total pending bytes to \(e\),
and changes latest-cut slack only by \(-N_{s,e}(j)\). Reject the proposed
change if any resulting latest-cut slack is negative, before making
any edits. These updates preserve the defining sum. To migrate across
distinct senders, check that the destination can take the row first;
if rejected, keep the old protected row and its ordering. If admitted,
append and retire the old row as one modeled atomic step. Distinct
sender budgets are essential to this elementary order of checks. The
new and old cuts must be at the actual queue owner for this model to
apply; no claim that observing a remote change makes it atomic.

\paragraph{Why only storing the latest projection fails.}
Two queues contain the same pair of 60-byte probes to \(E\) and \(F\)
and the same two 173-byte reports. One has \(E\) before the earlier report
and \(F\) between reports; the other swaps the probes. They agree on
each endpoint's total bytes and on every endpoint coefficient at the
latest protected cut. The same latest report is lawfully removed.
Both earlier reports now have the same 233-byte slack at capacity 466.
Growing \(E\) reduces the first queue's slack by 60 and leaves the
second unchanged. The fourth such edit must be refused in the first
case and can be accepted in the second. Thus the latest weighted
total, even together with the latest coefficient vector and the
rewound slack scalar, cannot support subsequent growth checks.
Saving order information or revisiting the queue is necessary in this
comparison. This is not a lower bound against every indexed data structure.

\paragraph{Cost and evidence.}
With hash tables at expected constant lookup cost and rank horizon \(T\),
one byte/rank update takes \(O(\log T)\); protected-heap pushes and
eventual purges add amortized \(O(\log T)\). A bind checks the \(d_e\)
current sender-endpoint incidences, \(O(d_e\log T)\) in the worst case
apart from amortized heap purges, instead of walking the pending queue.
Cancelling the latest and rewinding can read all \(d_s\) endpoints of its
old sender, \(O(d_s\log T)\). Migration charges both steps; pending rows,
sparse trees and stale buried heap keys cost \(O(n\log T+T)\) space
for \(n\) pending rows. No continual-horizon storage or real-time queue
lock bound follows. In 14,400 seeded steps, a full recomputation matched
all admissions and exact slacks, including 1,430 migration requests
(1,149 feasible). The indexed binding checks visited 4,236 incidences
and made 3,127 prefix queries; the scanning implementation visited
69,430 queue rows for binding. These are different primitive counts,
not a measured speed ratio; inspection, updates and 945 rewind-edge
visits are separately reported. Another direct solver with the same
index has the same cost, and a dense shared destination can erase the
advantage over scanning. Lean checks the indistinguishability pair.
There is no independent workload or general R2 acceptance.
\end{document}
